\documentclass[11pt]{amsart}
\usepackage[T1]{fontenc}
\usepackage[utf8]{inputenc}
\usepackage{amsmath,amssymb,amsthm,booktabs}
\usepackage[margin=1in]{geometry}
\usepackage[hidelinks]{hyperref}
\numberwithin{equation}{section}
\newtheorem{proposition}{Proposition}[section]
\theoremstyle{remark}
\newtheorem{observation}{Observation}
\newcommand{\norm}[1]{\left\lVert #1\right\rVert}
\newcommand{\abs}[1]{\left\lvert #1\right\rvert}
\newcommand{\Caputo}{{}^{C}_{0}D_t^\alpha}
\title{Verification notes on the L1/L1-2\texorpdfstring{\\}{ }Legendre spectral method}
\author{Hardik Bisht}
\date{November 2026}
\begin{document}
\maketitle

\begin{quote}
\small\textbf{Repository transcription and evidence notice (added at integration).}
This document was manually transcribed from the supplied eleven-page PDF
\texttt{legendre\_spectral\_verification\_notes\_st03.pdf}; the original
\LaTeX{} source was not supplied or recovered. Sections 1--7 retain the
source's logical order, hypotheses, numbered equations, propositions, and
qualifications. Page breaks are not preserved. Source-page comments and the
companion \texttt{PROVENANCE.md} locate the material.
Reports of executed experiments in the supplied notes are historical reports,
not claims of fresh execution during integration. Only the ST-03 Python script
and its text output were supplied alongside the PDF. Other companion programs
mentioned in Section 5.5 and the VQ-01 checker were not supplied or found.
The source-reported literature claims in Section 6.4 have not been rechecked
against the independent article. Actual integration checks and unavailable
steps are recorded separately in \texttt{results/VALIDATION.md}.
\end{quote}

% Supplied verification PDF, page 1.
\section{Purpose and reference problem}
This is a separate working document for the mathematical observations raised
during a study and numerical reproduction of the following article:
\begin{quote}
Y. A. Sabawi, S. Santra, P. Das, and N. A. Mohammad,
\emph{$H^1$ stability and convergence analysis of the L1/L1-2 Legendre
spectral method for non-local weakly singular integro-PDEs},
Calcolo \textbf{62}, article 42 (2025).
\href{https://doi.org/10.1007/s10092-025-00671-y}{doi:10.1007/s10092-025-00671-y}.
\end{quote}
The main model is a time-fractional initial--boundary value problem with a
weakly singular Volterra memory term. In one space dimension, its structure is
\begin{equation}\label{eq:model}
\Caputo u(x,t)-u_{xx}(x,t)+\lambda J[u](x,t)=f(x,t),
\qquad 0<\alpha,\beta<1,
\end{equation}
where
\begin{equation}\label{eq:memory}
J[u](x,t)=K(x,t)\int_0^t(t-s)^{-\beta}u(x,s)\,ds.
\end{equation}
The article uses L1 and L1-2 time discretizations, a combined
trapezoidal/singular-panel quadrature for the Volterra operator, and
Legendre--Gauss--Lobatto spectral discretization in space. Its analysis
discusses time-step well-posedness, stability, and temporal and spatial
convergence.

The purpose of this document is not to declare the reported theorems false.
It collects individual mathematical questions, the evidence currently
available, and the points that remain to be independently derived or tested.
An invalid step in a proof does not, by itself, imply that the theorem is
false. The original seven observations were initially provisional.
Observation 2 has a verified algebraic discrepancy (Section~\ref{sec:cf}),
while Observation 8 has a verified quadrature counterexample and sharp
consistency estimate (Section~\ref{sec:vq}). Section~\ref{sec:st} now
independently audits the L1--2 portion of the stability argument and identifies
a specific failure of its coefficient-based induction; the theorem itself
remains unresolved. The other observations retain provisional status.

\section{Observations transferred from the original writeup}
The following seven observations are reproduced from the appendix to the
main study notes. The entries record the initial reading; subsequent
verification is documented in the dedicated sections below.

\begin{observation}[Signed Volterra history terms in Theorem 3.1]
\textbf{Status: proof gap established at the coefficient-based induction
step; see Section~\ref{sec:st}.}
The energy estimate obtained by testing the weak equation with the current
numerical solution is a valid starting point. The subsequent step on page 16
bounds the norm of a signed sum of previous solutions by the absolute value
of the sum of the coefficients. In general,
% Supplied verification PDF, page 2.
\[
\norm{\sum_j c_jU^j}\leq\sum_j\abs{c_j}\norm{U^j},
\quad\text{but not necessarily}\quad
\norm{\sum_j c_jU^j}\leq\abs{\sum_j c_j}\max_j\norm{U^j}.
\]
This becomes consequential when the Volterra terms contribute negative
coefficients. The source's contraction proof therefore needs an additional
argument, stronger hypotheses, or a different stability estimate.
Cauchy--Schwarz itself is not the invalid step; the issue is how the signed
historical contributions are combined.
\end{observation}

\begin{observation}[L1-2 coefficients inferred from interpolation]
\textbf{Status: verified algebraic inconsistency in the published
coefficient representation; see Section~\ref{sec:cf}.}
From the L1-2 quadratic correction of equations (23)--(24),
\[
b_0^*=\frac1{2-\alpha}-\frac12=\frac{\alpha}{2(2-\alpha)}.
\]
Thus for steps after startup, direct collection of the forward differences gives
\[
d_0=b_0+b_0^*=1+\frac{\alpha}{2(2-\alpha)},
\]
which is not equal to one for $0<\alpha<1$. Equation (29), however, asserts
$d_0=1$ and bases later algebra on that identity. The index range for the
middle coefficient case also requires clarification. Preserve the two-sum
interpolatory L1-2 formula as an independent reference when implementing the
scheme.
\end{observation}

\begin{observation}[L1-2 startup and temporal consistency]
\textbf{Status: explicit test case available.}
The first panel still uses linear interpolation. For the smooth test function
$u(t)=t^2$, the quadratic interpolant is exact on the second panel but the
first-panel contribution to the Caputo error at $t_2=2\tau$ is
\[
\Phi_{L12}^{2}=\frac{\tau^{2-\alpha}}{\Gamma(1-\alpha)}
\int_0^1(2z-1)(2-z)^{-\alpha}\,dz.
\]
The integral is nonzero for $0<\alpha<1$. Hence this step has a contribution
of order $\tau^{2-\alpha}$, which is incompatible with a uniform
$O(\tau^{3-\alpha})$ Caputo truncation estimate for every step $n\geq2$
without a startup qualification. This should be checked separately from
convergence at fixed positive physical time.
\end{observation}

\begin{observation}[Quadrature consistency for variable coefficients]
\textbf{Status: exactness claim requires qualification.}
Theorem 4.1 invokes the equality $B(R_M,\rho_M)=B_M(R_M,\rho_M)$ when $R_M$
has a suitably restricted polynomial degree. Legendre--Gauss--Lobatto
quadrature with $M+1$ nodes is exact through degree $2M-1$, but the mass form
contains the spatially varying coefficient $\widehat R(x)$ and need not be a
polynomial of admissible degree. For $M=3$, take
\[
R_M=1-x^2,\qquad \rho_M=x(1-x^2),\qquad
\widehat R(x)=1+cx,\qquad 0<c<1.
\]
The exact integral and four-point Lobatto quadrature for the mass integrand
differ in absolute value by $32c/525$. A variable-coefficient quadrature
consistency term, or stronger hypotheses, appears necessary.
\end{observation}

\begin{observation}[Time-step-dependent $H^1$-type norm]
\textbf{Status: interpretation caveat.}
The norm used throughout the paper is
\[
\norm{v}_1^2=\norm{v}_{L^2}^2+\alpha_0\norm{v_x}_{L^2}^2,
\qquad \alpha_0=\Gamma(2-\alpha)\tau^\alpha.
\]
For a fixed $\tau>0$ it is equivalent to the standard $H^1$ norm, but the
equivalence constants are not uniform as $\tau\downarrow0$. Contractivity or
convergence in this weighted norm should not be silently restated as a
uniform estimate in the standard $H^1$ norm.
\end{observation}

% Supplied verification PDF, page 3.
\begin{observation}[Reported discrete error normalization]
\textbf{Status: numerical reproduction check.}
The numerical section defines discrete errors with spatial Lobatto weights
and an additional averaging factor involving $(M+1)^{-1}$. The factor is not
part of the standard Lobatto quadrature approximation of an $L^2$ norm,
because the quadrature weights already account for the spatial integration.
During reproduction, report the article's metric and the conventional
quadrature norm separately; otherwise apparent spatial rates may be affected
by the normalization.
\end{observation}

\begin{observation}[Kernel-sign assumption versus numerical examples]
\textbf{Status: assumptions check.}
The stability proof assumes a strictly positive spatial coefficient
$K(x,t)$. The paper's numerical Example 6.2 uses a factor
$K=x$ on a domain containing negative $x$, so that hypothesis is
not satisfied throughout the domain. Distinguish which results apply under
positivity and which experiments probe more general coefficients.
\end{observation}

\section{Additional observation identified subsequently}
\begin{observation}[Order of the weakly singular Volterra quadrature]
\textbf{Status: verified for equation (6) as printed; see
Section~\ref{sec:vq}.}
Equation (6) applies ordinary composite trapezoidal quadrature on all but the
final time panel, using a special singular-panel rule on the final interval.
The paper's equation (8) reports a Volterra quadrature remainder of order
$O(\tau^{2-\beta})$ under its stated smoothness regime. The ordinary
trapezoidal contributions near the weak singularity require closer scrutiny.

A simple test is $K=1$, $u(s)\equiv1$, $T=1$, and $\beta=1/2$.
The exact integral is
\[
\int_0^1(1-s)^{-1/2}\,ds=2.
\]
For the ordinary trapezoidal panel immediately preceding the final special
panel, $[1-2\tau,1-\tau]$, the difference between the trapezoidal
approximation and the exact integral is
\[
\left[\frac{1+2^{-1/2}}2-2(\sqrt2-1)\right]\tau^{1/2}.
\]
The bracket is strictly positive. Since the kernel is convex on earlier
panels and the special final-panel contribution is exact for constant $u$,
the combined error cannot have a uniform $O(\tau^{3/2})$ bound in this test.
This contradicts the stated uniform bound. Section~\ref{sec:vq} supplies a
rigorous proof of the sharp general $O(\tau^{1-\beta})$ order and
independently evaluated numerical data. The effect on the full PDE solution
error remains an open task.
\end{observation}

\section{Verification VQ-01: the printed Volterra quadrature}\label{sec:vq}
\subsection{Claim and assumptions}
The article's equation (6), on PDF page 7 (journal section 2.1), treats every
time panel except the last with the ordinary composite trapezoidal rule; the
final, singular panel uses an endpoint average multiplied by the analytically
integrated kernel weight. Equation (8) states that the quadrature error is
$O((\Delta t)^{2-\beta})$. We check exactly this printed construction, not a
product-integration method or a differently weighted trapezoidal rule.

Fix a spatial point, set $t_n=n\tau$, $\tau=T/N$, and suppress the
$x$-argument. For the separable kernel
$K(t,s)=K(t)(t-s)^{-\beta}$, define
\begin{equation}\label{eq:vq-integral}
I^n[u]=K(t_n)\int_0^{t_n}(t_n-s)^{-\beta}u(s)\,ds.
\end{equation}
% Supplied verification PDF, page 4.
The numerical formula in equation (6) is
\begin{equation}\label{eq:vq-rule}
\begin{split}
Q_\tau^n[u]=K(t_n)\biggl\{
&\frac\tau2\sum_{j=0}^{n-2}
\bigl[(t_n-t_j)^{-\beta}u(t_j)+(t_n-t_{j+1})^{-\beta}u(t_{j+1})\bigr]\\
&+\frac{\tau^{1-\beta}}{2(1-\beta)}[u(t_{n-1})+u(t_n)]\biggr\}.
\end{split}
\end{equation}
We denote the quadrature error by $\Phi_T^n=I^n[u]-Q_\tau^n[u]$. This sign
convention agrees with the article's definition, though the absolute-error
conclusions below are independent of that sign.

\subsection{A counterexample for every $0<\beta<1$}
\begin{proposition}[Failure of the claimed quadrature order]\label{prop:vq-lower}
Let $T=1$, $K\equiv1$, $u\equiv1$, $0<\beta<1$, and $n=N\geq2$.
Then
\begin{equation}\label{eq:vq-lower}
\abs{\Phi_T^N}\geq c_\beta\tau^{1-\beta},\qquad
c_\beta:=\frac{1+2^{-\beta}}2-\frac{2^{1-\beta}-1}{1-\beta}>0.
\end{equation}
In particular, there exists no $\tau$-independent constant $C$ such that
$\abs{\Phi_T^N}\leq C\tau^{2-\beta}$ for all sufficiently small $\tau$.
\end{proposition}
\begin{proof}
The exact value is
\[
I^N[1]=\int_0^1(1-s)^{-\beta}\,ds=\frac1{1-\beta}.
\]
The special last-panel term in \eqref{eq:vq-rule} equals
$\tau^{1-\beta}/(1-\beta)$ and therefore integrates its panel exactly when
$u=1$. On the immediately preceding panel $[1-2\tau,1-\tau]$, the ordinary
trapezoidal approximation equals
\[
Q_{\mathrm{near}}=\frac{\tau^{1-\beta}}2(1+2^{-\beta}),
\]
whereas the exact panel integral is
\[
I_{\mathrm{near}}=\tau^{1-\beta}\frac{2^{1-\beta}-1}{1-\beta}.
\]
Their difference is exactly $c_\beta\tau^{1-\beta}$. The function
$s\mapsto(1-s)^{-\beta}$ has strictly positive second derivative wherever
$s<1$:
\[
\frac{d^2}{ds^2}(1-s)^{-\beta}=\beta(\beta+1)(1-s)^{-\beta-2}>0.
\]
Therefore every ordinary trapezoidal panel has a strictly positive
overestimate; none can cancel the near-panel error. This proves
\eqref{eq:vq-lower}. Dividing by $\tau^{2-\beta}$ gives
$\abs{\Phi_T^N}/\tau^{2-\beta}\geq c_\beta/\tau\to\infty$, contradicting
any uniform bound of the reported order.
\end{proof}
In fact, on this constant-function test, we have the exact identity
\begin{equation}\label{eq:vq-series}
Q_\tau^N[1]-I^N[1]=\tau^{1-\beta}\sum_{r=2}^N
\left[\frac{(r-1)^{-\beta}+r^{-\beta}}2
-\frac{r^{1-\beta}-(r-1)^{1-\beta}}{1-\beta}\right].
\end{equation}
Each bracket is positive. As $r\to\infty$ it is asymptotic to
$\beta(\beta+1)r^{-\beta-2}/12$, so the series converges to a finite positive
number $C_\beta$. Consequently,
\begin{equation}\label{eq:vq-asymptotic}
\abs{\Phi_T^N}\sim C_\beta\tau^{1-\beta}\qquad(N\to\infty,\ T=1).
\end{equation}
Thus the loss of one power is not merely an artifact of a lower bound.

% Supplied verification PDF, page 5.
\subsection{General upper estimate}
\begin{proposition}[Sharp general-order estimate]\label{prop:vq-upper}
Assume $u\in C^2([0,T])$, $0<\beta<1$, and $K$ is bounded on
$[0,T]$. Then the printed formula satisfies
\[
\max_{1\leq n\leq N}\abs{I^n[u]-Q_\tau^n[u]}\leq C\tau^{1-\beta},
\]
where $C$ depends on $T,\beta,\norm{K}_\infty$ and the indicated
bounds on $u$, but not on $N$. The power $1-\beta$ cannot be improved
uniformly for all such functions.
\end{proposition}
\begin{proof}
Fix $t_n$ and put $g_n(s)=(t_n-s)^{-\beta}u(s)$. Direct differentiation gives
\begin{equation}\label{eq:vq-gpp}
\begin{split}
g_n''(s)={}&\beta(\beta+1)(t_n-s)^{-\beta-2}u(s)
 +2\beta(t_n-s)^{-\beta-1}u'(s)\\
&+(t_n-s)^{-\beta}u''(s).
\end{split}
\end{equation}
Index the ordinary panels backwards from $t_n$, so panel $r$ is
$[t_n-r\tau,t_n-(r-1)\tau]$, $r=2,\ldots,n$. On this panel,
$t_n-s\geq(r-1)\tau$. The standard local trapezoidal-error estimate yields
\[
\begin{split}
\abs{E_r}\leq C\bigl[&\norm{u}_\infty\tau^{1-\beta}(r-1)^{-\beta-2}
 +\norm{u'}_\infty\tau^{2-\beta}(r-1)^{-\beta-1}\\
&+\norm{u''}_\infty\tau^{3-\beta}(r-1)^{-\beta}\bigr].
\end{split}
\]
Both $\sum_{r=1}^{\infty}r^{-\beta-2}$ and
$\sum_{r=1}^{\infty}r^{-\beta-1}$ are finite, while
$\sum_{r=1}^{n-1}r^{-\beta}\leq C_\beta n^{1-\beta}$. Since $n\tau\leq T$,
summing the ordinary-panel errors gives
\[
\sum_{r=2}^n\abs{E_r}\leq C(\tau^{1-\beta}+\tau^{2-\beta}+\tau^2).
\]
For the final panel, the difference between $u(s)$ and the average of its
endpoint values is at most $\tau\norm{u'}_\infty/2$. Integrating the positive
weight gives an error bounded by $C\tau^{2-\beta}$. Multiply by
$\norm{K}_\infty$ and combine the panels. For fixed $T$ and
sufficiently small $\tau$, the bound is $C\tau^{1-\beta}$. Sharpness follows
from Proposition~\ref{prop:vq-lower}.
\end{proof}

\subsection{Independent numerical check}
\textit{Historical numerical evidence reported in the supplied notes.}
For the exact test $u=K=1$, $T=1$, we evaluated the finite sum in
\eqref{eq:vq-rule} and compared it with $I^N=1/(1-\beta)$. The empirical
order uses $p_N=\log_2(E_N/E_{2N})$.
\begin{center}
\begin{tabular}{@{}ccccc@{}}
\toprule
$\beta$ & $E_{256}$ & $E_{512}$ & $E_{1024}$ & $p_{512}$\\
\midrule
$0.2$ & $1.90147\cdot10^{-4}$ & $1.09293\cdot10^{-4}$ & $6.27930\cdot10^{-5}$ & $0.7995$\\
$0.5$ & $2.47721\cdot10^{-3}$ & $1.75194\cdot10^{-3}$ & $1.23888\cdot10^{-3}$ & $0.4999$\\
$0.8$ & $2.06036\cdot10^{-2}$ & $1.79371\cdot10^{-2}$ & $1.56153\cdot10^{-2}$ & $0.2000$\\
\bottomrule
\end{tabular}
\end{center}
In each case $p_{512}$ approaches $1-\beta$, rather than the paper's
$2-\beta$. Two equivalent finite-sum implementations, one using the original
time nodes and one using \eqref{eq:vq-series}, agree to floating-point
precision. A companion Python checker reproduces the table and its estimated
orders; no PDE solver is required.
\textit{Integration qualification: that companion checker was not supplied;
the new regression checker is separately identified in the repository.}

\subsection{Scope and consequences}
The constant-in-time test is admissible at the level of the quadrature
estimate. It can also be embedded in a homogeneous Dirichlet IBVP: set
$\Omega_x=(-1,1)$,
\[
u(x,t)=1-x^2,\qquad K(x,t)=1,\qquad
f(x,t)=2+\lambda(1-x^2)\frac{t^{1-\beta}}{1-\beta},
\]
with initial datum $1-x^2$ and zero boundary values. The Caputo derivative
vanishes, and at every interior point with $u(x,t)\neq0$, the same Volterra
error persists after multiplication by $1-x^2$.

With the conventional definition of the operator discrepancy,
\[
r_\tau^n=\Phi_{L1}^n+\lambda\Phi_T^n,
\]
the general operator consistency estimate from the printed rules is at most
\[
\abs{r_\tau^n}\leq C(\tau^{2-\alpha}+\tau^{1-\beta})=O(\tau^{1-\beta})
\]
% Supplied verification PDF, page 6.
under standard smoothness assumptions. On the constant-in-time example,
$\Phi_{L1}^n=0$, so the lower order occurs in the complete operator
discrepancy whenever $\lambda\neq0$. The article's scaled residual convention
and signs should be checked independently; multiplying this Volterra
discrepancy by $\alpha_0=\Gamma(2-\alpha)\tau^\alpha$ yields a term of order
$\tau^{1+\alpha-\beta}$.

\textbf{What is established:} equation (6), as printed, cannot have the
reported general $O(\tau^{2-\beta})$ Volterra quadrature remainder in equation
(8). Its sharp general consistency order for smooth $u$ is $1-\beta$.

\textbf{What is not established:} this does not by itself disprove the paper's
global PDE solution-error estimates, nor determine which corrected stability
and convergence bounds hold. Those require a separate analysis of the
discrete evolution and, where appropriate, experiments with the full scheme.

\subsection{A candidate higher-order replacement}
One possible alternative is product integration: interpolate $u$ piecewise
linearly, but integrate the singular weight exactly on each time panel. For
the interpolant $I_\tau u$,
\[
Q_{\mathrm{PI}}^n[u]=K(t_n)\sum_{j=0}^{n-1}
\int_{t_j}^{t_{j+1}}(t_n-s)^{-\beta}(I_\tau u)(s)\,ds.
\]
Since $\norm{u-I_\tau u}_\infty\leq\tau^2\norm{u''}_\infty/8$, this
construction satisfies
\[
\abs{I^n[u]-Q_{\mathrm{PI}}^n[u]}
\leq\frac{\norm{K}_\infty T^{1-\beta}}{8(1-\beta)}
\norm{u''}_\infty\tau^2.
\]
This is not the method printed in equation (6). Its discrete weights,
stability, and full-scheme convergence must be investigated separately before
it can be recommended as a replacement.

\section{Verification CF-02: L1--2 coefficient identities}\label{sec:cf}
\subsection{Source statement and what is being checked}
In Section 2.2 of the reference article, equations (20)--(22) define the
L1--2 approximation by linearly interpolating on the first time panel and
using a backward three-point quadratic interpolant thereafter. Equation (23),
on PDF page 11, first gives a two-sum formula, then claims that it equals a
single sum of backward increments weighted by $d_j$. Equation (24) defines
the correction weights; equation (29), on PDF page 13, states that $d_0=1$
and that the $d_j$ are positive and decreasing. The recurrence in equations
(28)--(32) uses these assertions.

We test the algebraic equivalence of the formulas as printed, independently
of the spatial problem, the Volterra quadrature, and the PDE stability
theorem. A discrepancy in an identity is not, on its own, evidence that the
interpolation-based scheme is unstable.

Write $\tau=\Delta t$, $u^j=u(x,t_j)$ at a fixed spatial point, and
$\Delta u_{j}=u^{j+1}-u^j$. The first nontrivial L1--2 step occurs at $t_2$
($k=1$). The two-sum formula directly following from the defining
interpolants is
\begin{equation}\label{eq:cf-two}
D_{L12}^{\alpha}u^{k+1}=\frac{\tau^{-\alpha}}{\Gamma(2-\alpha)}
\left[\sum_{j=0}^kb_{k-j}\Delta u_{j}+
\sum_{j=1}^kb_{k-j}^*(\Delta u_{j}-\Delta u_{j-1})\right],
\end{equation}
where $b_r=(r+1)^{1-\alpha}-r^{1-\alpha}$ and
\begin{equation}\label{eq:cf-star}
b_r^*=\frac{(r+1)^{2-\alpha}-r^{2-\alpha}}{2-\alpha}
-\frac{(r+1)^{1-\alpha}+r^{1-\alpha}}2.
\end{equation}
In fact, using the local coordinate $z=(s-t_j)/\tau$, the quadratic
correction comes from the moment
\[
b_r^*=(1-\alpha)\int_0^1(z-\tfrac12)(r+1-z)^{-\alpha}\,dz.
\]
Evaluating this integral yields \eqref{eq:cf-star}, independently of the
single-sum representation in the article.

% Supplied verification PDF, page 7.
\subsection{Exact collection of the forward differences}
\begin{proposition}[Correct coefficients for the printed interpolatory rule]\label{prop:cf-collect}
For the L1--2 discretization \eqref{eq:cf-two}, the correct collected
representation is
\begin{equation}\label{eq:cf-one}
D_{L12}^{\alpha}u^{k+1}=\frac{\tau^{-\alpha}}{\Gamma(2-\alpha)}
\sum_{r=0}^kd_r^{[k]}\Delta u_{k-r},
\end{equation}
with startup $d_0^{[0]}=1$ and, when $k\geq1$,
\begin{equation}\label{eq:cf-coeff}
d_r^{[k]}=\begin{cases}
b_0+b_0^*, & r=0,\\
b_r+b_r^*-b_{r-1}^*, & 1\leq r\leq k-1,\\
b_k-b_{k-1}^*, & r=k.
\end{cases}
\end{equation}
The middle case is empty when $k=1$.
\end{proposition}
\begin{proof}
In the second sum of \eqref{eq:cf-two}, each term equals
$b_{k-j}^*\Delta u_{j}-b_{k-j}^*\Delta u_{j-1}$. The newest increment
$\Delta u_{k}$ receives $b_0$ from the first sum and $b_0^*$ from the second.
Each interior increment $\Delta u_{k-r}$, $1\leq r\leq k-1$, receives
$b_r+b_r^*-b_{r-1}^*$. The oldest increment $\Delta u_{0}$ receives
$b_k-b_{k-1}^*$. For $k=0$, the second sum is empty and $b_0=1$. These are
exactly the coefficients in \eqref{eq:cf-coeff}.
\end{proof}
The superscript $[k]$ is intentional: the oldest-panel coefficient depends
on which step is being computed. A single infinite sequence $d_0,d_1,\ldots$
does not encode that endpoint modification without specifying the time level.

\subsection{Two exact counterexamples to the printed identities}
\begin{proposition}[The leading coefficient and positivity claim]\label{prop:cf-counter}
For every $0<\alpha<1$ and every $k\geq1$, the coefficient of $\Delta u_{k}$
in \eqref{eq:cf-two} satisfies
\begin{equation}\label{eq:cf-leading}
d_0^{[k]}=1+b_0^*=1+\frac{\alpha}{2(2-\alpha)}>1.
\end{equation}
Moreover, for the admissible value $\alpha=3/4$ at $k=1$,
\begin{equation}\label{eq:cf-negative}
d_0^{[1]}=\frac{13}{10},\qquad d_1^{[1]}=2^{1/4}-\frac{13}{10}<0.
\end{equation}
Thus neither $d_0=1$ for all steps nor $d_j>0$ for all $0<\alpha<1$ follows
from equation (23).
\end{proposition}
\begin{proof}
Putting $r=0$ in \eqref{eq:cf-star} gives
$b_0^*=1/(2-\alpha)-1/2=\alpha/[2(2-\alpha)]$, which is strictly positive.
For $k=1$, the oldest-increment coefficient is $d_1^{[1]}=b_1-b_0^*$.
At $\alpha=3/4$, this is $2^{1/4}-13/10$; it is negative since
$2<(13/10)^4=28561/10000$, and both quantities are positive.
\end{proof}
There is an additional, independent index reversal in the final line of the
article's equation (23). It displays a sum of the form
\[
\sum_{j=0}^k\Delta u_{k-j}d_{k-j},
\]
whereas the two-sum interpolation gives
$\sum_{j=0}^k\Delta u_{k-j}d_j^{[k]}$, equivalently
$\sum_{j=0}^k\Delta u_{j}d_{k-j}^{[k]}$. For example, at $k=1$ choose nodal
data $(u^0,u^1,u^2)=(0,0,1)$. The two-sum formula is proportional to
$d_0^{[1]}=13/10$ for $\alpha=3/4$, while the final single sum as printed is
proportional to $d_1^{[1]}=2^{1/4}-13/10<0$. The two expressions are plainly
not equal.

Finally, equation (24) prints the index range $j=0,2,\ldots,N-1$ for
$b_j^*$, apparently omitting $j=1$; the middle case of the displayed
definition of $d_j$ also begins at $j=2$ instead of $j=1$. Both indices are
needed for a correct representation when $k\geq2$.

% Supplied verification PDF, page 8.
\subsection{Consequences for the recurrence and stability arguments}
For any fixed step $k\geq1$, the basic telescoping identity is
\begin{equation}\label{eq:cf-telescoping}
(d_0^{[k]}-d_1^{[k]})+\sum_{r=1}^{k-1}(d_r^{[k]}-d_{r+1}^{[k]})
+d_k^{[k]}=d_0^{[k]},
\end{equation}
not 1. This conflicts with the normalization asserted in equation (29).
More strongly, the corrected $d_1^{[1]}<0$ at $\alpha=3/4$, so positivity
and monotonicity-based inductions cannot be justified using (29) without
additional analysis.

In the elliptic rearrangement of the interpolatory L1--2 scheme, the
current-state coefficient should accordingly be
\[
\widehat R_{(2)}^{k+1}(x)=d_0^{[k]}
+\frac{\lambda\alpha_0\tau^{1-\beta}}{2(1-\beta)}K(x,t_{k+1}),
\qquad \alpha_0=\Gamma(2-\alpha)\tau^\alpha,
\]
for $k\geq1$, rather than replacing $d_0^{[k]}$ by 1 as in equation (31).
It is possible to normalize the entire equation by $d_0^{[k]}$, but that
would also change the diffusion coefficient to $\alpha_0/d_0^{[k]}$ and
rescale all historical coefficients; merely writing $d_0=1$ is not such a
normalization. The startup step $k=0$ is unaffected.

This is an algebraic verification of the coefficient identities, not a
stability counterexample. The implications for the numerical scheme and
Theorem 3.1 remain open. In particular, a negative historical coefficient
does not by itself imply numerical instability.

\subsection{Independent reproducibility and formalization status}
\textit{Historical execution and companion-file report in the supplied notes.}
We separately evaluated the defining piecewise-interpolatory Caputo integrals
numerically, the two-sum formula \eqref{eq:cf-two}, and the corrected single
sum \eqref{eq:cf-one}, for $\alpha\in\{0.25,0.5,0.75\}$,
$k\in\{0,1,2,4\}$, and polynomial nodal data of degrees 0 through 3.
The two algebraic representations agree to high precision; direct panel
integration also agrees within numerical integration tolerance for the
tested steps. The last line of equation (23), used literally, fails the
simple $k=1$ test above.
\begin{center}
\begin{tabular}{@{}cccc@{}}
\toprule
$\alpha$ & $d_0^{[1]}$ & $d_1^{[1]}$ & $d_1^{[1]}>0$?\\
\midrule
$0.25$ & $1.07142857$ & $0.61036426$ & yes\\
$0.50$ & $1.16666667$ & $0.24754690$ & yes\\
$0.75$ & $1.30000000$ & $-0.11079288$ & no\\
\bottomrule
\end{tabular}
\end{center}
The accompanying Python script has been executed and passes the checks.
A standard-library Julia counterpart and a small mathlib Lean draft are also
provided. Neither the Julia file nor the Lean proof draft was executed in
this environment; hence this document does not claim machine-certified
formal verification. The exact algebraic proofs above do not depend on those
programs.
\textit{Integration qualification: the CF-02 Python/Julia/Lean companions
mentioned in this historical paragraph are absent from the supplied files
and repository inventory; their claimed presence and execution cannot be
independently confirmed here. The supplied ST-03 script is a different file.}

\paragraph{Status.}
The published single-sum coefficient representation and equation (29) are
algebraically inconsistent with the defining two-sum interpolation for the
stated parameter range. The correct interpolation-based scheme has been
reconstructed. Whether the paper's implemented code follows the two-sum rule
or the misindexed recurrence, and whether its reported stability conclusions
can be recovered, remain open.

\section{Verification ST-03: L1--2 history coefficients and Theorem 3.1}\label{sec:st}
\subsection{Scope, source claim, and reduced problem}
Theorem 3.1 (Section 3 of Sabawi et al.) claims, subject to a sufficiently
fine temporal mesh, the bound
\begin{equation}\label{eq:st-claim}
\norm{U_{(q)}^{k+1}}_1\leq\norm{U_{(q)}^0}_{L^2},\qquad q=1,2,
\end{equation}
where $\norm{v}_1^2=\norm{v}_{L^2}^2+\alpha_0\norm{v_x}_{L^2}^2$ and
$\alpha_0=\Gamma(2-\alpha)\tau^\alpha$. Its written proof establishes a
first-step energy estimate and then performs induction for $q=1$, asserting
that the $q=2$ case follows similarly from equation (29). Our objective here
is not to prove or disprove \eqref{eq:st-claim}, but to determine whether
that induction applies to the interpolatory L1--2 scheme actually defined
in equation (23).

% Supplied verification PDF, page 9.
For the stability test set the forcing to zero and work at a fixed time
level $t_{k+1}$; write $U^j=U_{(2)}^j$ and $\tau=\Delta t$. The following
derivation first isolates the fractional derivative and then reincorporates
the Volterra history term. We assume real-valued solutions, as in the
article's energy formulation.

\subsection{The correct state-value recurrence}
Let $d_r^{[k]}$ denote the interpolatory coefficients independently derived
in Proposition~\ref{prop:cf-collect}, and define the historical coefficients
\begin{equation}\label{eq:st-history}
h_r^{[k]}=d_r^{[k]}-d_{r+1}^{[k]},\quad 0\leq r<k,
\qquad h_k^{[k]}=d_k^{[k]}.
\end{equation}
The backward-increment representation of the fractional operator gives the
exact algebraic identity
\begin{equation}\label{eq:st-state}
\sum_{r=0}^kd_r^{[k]}(U^{k+1-r}-U^{k-r})
=d_0^{[k]}U^{k+1}-\sum_{r=0}^kh_r^{[k]}U^{k-r}.
\end{equation}
The history coefficients telescope to
\begin{equation}\label{eq:st-telescope}
\sum_{r=0}^kh_r^{[k]}=d_0^{[k]},
\end{equation}
which equals 1 only at startup. Identity \eqref{eq:st-telescope} does not
imply that the coefficients $h_r^{[k]}$ are nonnegative, nor that their
absolute values sum to $d_0^{[k]}$.

For a spatially constant positive coefficient $K=1$, define
\[
\theta=\frac{\lambda\alpha_0\tau^{1-\beta}}{2(1-\beta)},\qquad
Q_{\mathrm{past}}^{k+1}=\frac\tau2\sum_{j=0}^{k-1}
\bigl[((k+1-j)\tau)^{-\beta}U^j+((k-j)\tau)^{-\beta}U^{j+1}\bigr].
\]
Multiplying the homogeneous time-discrete equation by $\alpha_0$, the
corrected recurrence becomes
\begin{equation}\label{eq:st-recurrence}
(d_0^{[k]}+\theta)U^{k+1}-\alpha_0U_{xx}^{k+1}
=\sum_{r=0}^kh_r^{[k]}U^{k-r}-\theta U^k
-\lambda\alpha_0Q_{\mathrm{past}}^{k+1}\qquad(k\geq1).
\end{equation}
If $K=K(x,t)$ varies spatially, the $\theta$ terms are
multiplied by $K(x,t_{k+1})$ and the past-panel terms by the
corresponding kernel factors; the key signed-history issue remains.
The current-state coefficient in \eqref{eq:st-recurrence} is
$d_0^{[k]}+\theta$, not $1+\theta$, as printed in equations (30)--(31).
At $k=0$, $d_0^{[0]}=1$ and the paper's startup coefficient is consistent.

\subsection{A certified failure of the required induction inequality}
\begin{proposition}[The elementary contraction induction cannot be justified
from the stated hypotheses]\label{prop:st-gap}
For the interpolatory L1--2 scheme with $\alpha=3/4$ and $k=1$, the
historical coefficients satisfy
\[
d_0^{[1]}=\frac{13}{10},\quad
d_1^{[1]}=2^{1/4}-\frac{13}{10}<0,\quad
h_0^{[1]}=d_0^{[1]}-d_1^{[1]},\quad h_1^{[1]}=d_1^{[1]}.
\]
In particular,
\begin{equation}\label{eq:st-budget}
\abs{h_0^{[1]}}+\abs{h_1^{[1]}}=d_0^{[1]}-2d_1^{[1]}>d_0^{[1]}.
\end{equation}
Hence an induction hypothesis consisting only of
$\norm{U^j}\leq\norm{U^0}$ cannot imply
$\norm{\sum_{r=0}^kh_r^{[k]}U^{k-r}}\leq d_0^{[k]}\norm{U^0}$ for all
possible preceding states.
\end{proposition}
\begin{proof}
The coefficient calculation is exact by Proposition~\ref{prop:cf-counter}.
For arbitrary real scalar states with $U^1=1$ and $U^0=-1$ (both having
norm one), the history combination is
$h_0^{[1]}-h_1^{[1]}=d_0^{[1]}-2d_1^{[1]}>d_0^{[1]}$. This refutes the
inference from norm bounds alone. These two arbitrarily assigned states are
not asserted to satisfy the fractional evolution recurrence, so this is a
proof-step counterexample, not a theorem counterexample.
\end{proof}

% Supplied verification PDF, page 10.
Numerically, $d_0^{[1]}=1.300000$, $d_1^{[1]}=-0.110792885$, and the
absolute historical coefficient budget is $1.521585770$, approximately
$1.170451$ times the current coefficient. Moreover, a negative historical
coefficient can arise even without $d_1<0$: at $\alpha=1/2$ and $k=2$,
$d_2^{[2]}>d_1^{[2]}$, so $h_1^{[2]}=d_1^{[2]}-d_2^{[2]}<0$. Thus the
naive nonnegative-coefficient induction is not generally available even for
moderate $\alpha$.

The Volterra contributions do not repair this inference automatically.
For a sum of signed history terms one must bound
\[
\norm{\sum_jc_jU^j}\leq\sum_j\abs{c_j}\norm{U^j},
\]
not replace $\sum_j\abs{c_j}$ with $\abs{\sum_jc_j}$.
Positivity of the kernel $K$ alone does not supply the missing
sign information on the combined coefficients. The problematic replacement
occurs explicitly in the published $q=1$ argument as well; the $q=2$ proof
is asserted by analogy rather than worked out.

\subsection{Independent literature: the exceptional coefficients are known}
\textit{Literature interpretation retained from the supplied notes; the
independent source was not available locally for rechecking at integration.}
A directly comparable analysis is M. Ramezani, R. Mokhtari, and G. Haase,
\emph{Analysis of stability and convergence for L-type formulas combined
with a spatial finite element method for solving subdiffusion problems},
Electronic Transactions on Numerical Analysis \textbf{55} (2022), 568--584,
\href{https://doi.org/10.1553/etna_vol55s568}{doi:10.1553/etna\_vol55s568}.

Section 3.2 (Theorems 3.4--3.6) treats L1--2 stability and convergence using
energy estimates and a discrete Gronwall argument. The authors expressly
note that strict coefficient monotonicity can fail for higher-order L-type
schemes, making elementary induction inappropriate.

This is independent corroboration that the exceptional coefficient structure
must be handled. But their model is a subdiffusion equation without Sabawi
et al.'s additional weakly singular Volterra memory term, and their
$L^2/H^1$ results cannot be imported as a proof of \eqref{eq:st-claim}.
In particular, we must first derive a discrete estimate for the combined
fractional and Volterra histories before invoking any Gronwall inequality.

\subsection{A reduced evolution test: why signed coefficients do not prove instability}
Consider a single diffusion eigenmode, with $\lambda=0$ solely to isolate
the L1--2 derivative and $-U_{xx}=\sigma U$, $\sigma\geq0$.
Put $\mu=\sigma\alpha_0\geq0$ and normalize the initial amplitude to
$u^0=1$. The first step gives $u^1=(1+\mu)^{-1}$. The second step, with the
correct L1--2 coefficients, satisfies
\begin{equation}\label{eq:st-mode}
(d_0+\mu)u^2=(d_0-d_1)u^1+d_1u^0
\quad\Longrightarrow\quad
u^2=\frac{d_0+d_1\mu}{(d_0+\mu)(1+\mu)}.
\end{equation}
For $0<\alpha<1$ we have $\abs{d_1}<1$ and $d_0>0$, hence
\[
\abs{u^2}\leq\frac{d_0+\abs{d_1}\mu}{(d_0+\mu)(1+\mu)}
\leq\frac1{1+\mu}.
\]
It follows in this particular scalar eigenmode that the modified norm
factor $\sqrt{1+\mu}\abs{u^2}$ is at most $(1+\mu)^{-1/2}\leq1$.
Thus the invalid induction step is compatible with genuine contractivity
of a reduced, exactly solved trajectory. This is only a two-step calculation
and does not extend to the full space-dependent Volterra scheme without
further proof.

The printed equation (30), with $d_0$ replaced by 1 but historical terms
left unscaled, defines a different recurrence. For example, at $\alpha=3/4$,
$\mu=0.3$, and zero memory coupling, the correct rule gives
$u^2\approx0.609020257$, whereas the unrescaled $d_0=1$ form gives
$u^2\approx0.572048600$. This diagnostic uses $\lambda=0$ to expose the
algebraic difference, not to assert a counterexample under the article's
$\lambda>0$ assumptions.

% Supplied verification PDF, page 11.
\subsection{Reproducible numerical diagnostic and present verdict}
\textit{Historical diagnostic reported in the notes and supplied text output.}
The accompanying \texttt{verify\_l12\_stability.py} independently evaluates
the collected state-value recurrence and the direct backward-increment
recurrence for the scalar equation
\[
\Caputo u+\sigma u+\lambda\int_0^t(t-s)^{-\beta}u(s)\,ds=0,
\qquad u(0)=1,\quad K=1.
\]
Across 432 parameter combinations, the two independently programmed
recurrence forms agree to within $4\times10^{-15}$ (ordinary double
precision). This checks our algebraic implementation, not the full PDE.
Some coarse-step tests exhibit weighted amplification above one; for example
$\alpha=0.25$, $\beta=0.8$, $\sigma=10$, $\lambda=10$, $N=16$ gives
$\sqrt{1+\sigma\alpha_0}\abs{u^1}\approx1.086$. This does not contradict
the article's sufficiently-small-step formulation, and here $\beta>\alpha$.
For fixed parameters, refinement in that test reduces the amplification
below one. No counterexample to the stated fine-grid theorem is claimed.

\paragraph{Verdict (ST-03).}
The L1--2 extension of the published Theorem 3.1 induction is not justified
by the stated coefficient properties: the normalization $d_0=1$ after
startup is incorrect, and signed history coefficients violate the required
absolute-sum bound. The independent 2022 L1--2 analysis explicitly recognizes
this sign difficulty and develops a different route to stability. The
stability theorem for Sabawi et al.'s full Volterra--diffusion scheme
remains unresolved; a corrected energy estimate with a suitable discrete
memory/Gronwall argument, or a parameter-valid counterexample, is still needed.

\section{Verification protocol}
Each observation will be considered individually. For each one, the
verification record should identify: (i) the exact statement and assumptions
in the article; (ii) an independent derivation or counterexample; (iii) a
numerical experiment when relevant; and (iv) whether the issue is a notation
error, an unjustified proof step, a counterexample to a stated estimate, or
an unresolved question. Mathematical conclusions should be updated only
when the evidence is sufficient.

\paragraph{Proposed starting point.}
Observations 2 and 8 have verified formula-level discrepancies, recorded in
Sections~\ref{sec:cf} and~\ref{sec:vq}. The L1--2 induction audit ST-03
(Section~\ref{sec:st}) verifies a proof-level gap, but not a counterexample
to Theorem 3.1. Next derive a stability estimate for the full
Volterra--diffusion recurrence using signed-coefficient bounds and then
reassess the claimed solution convergence rates.
\end{document}
